\chapter{Heat, Radiation, and Resilient Infrastructure}
\label{ch:22}

The terrestrial programs described in the preceding chapters share three physical constraints that are easily overlooked in the enthusiasm for each, and the present chapter, which closes the fourth part, treats them together. The first is heat: every use of energy ends as heat, and a civilization that multiplies its energy use will eventually meet the limit set by the planet's capacity to radiate it. The second is radiation, in the ordinary sense of ionizing radiation from nuclear processes and in the sense of the radiation environment that a long-lived structure experiences. The third is resilience, the capacity of an infrastructure to withstand shocks and to restore itself, which is the property that determines whether the programs deserve the label of infrastructure or are merely demonstrations. The three share a mathematical character, since in each there is a simple model whose parameters are measured and whose conclusions are sharper than the discussion usually allows. The mitigation pathways and the infrastructure constraints against which the heat budget is read are those of the sixth assessment of the intergovernmental panel~\cite{ipccwg3}.

The planetary heat budget supplies the first example. The Earth absorbs from the Sun about $1.2\times10^{17}$ W, which is a quarter of the solar constant times one minus the albedo times the surface area, and radiates the same amount to space in equilibrium. Human use of primary energy is about $1.8\times10^{13}$ W, which is $0.015$ per cent of the absorbed flux, or about $0.036\ \mathrm{W\,m^{-2}}$ averaged over the surface. The direct heating is accordingly negligible by comparison with the radiative forcing from greenhouse gases, which exceeds two watts per square meter, some sixty times larger. The comparison reassures about the present and misleads about the future, because the heat of human activity scales with the energy used whereas the capacity of the planet to shed it is fixed, and the limit would be reached by a growth of some orders of magnitude. If the climate responds to a forcing $F$ with an equilibrium warming $F/\lambda$, where the feedback parameter $\lambda$ is of the order of $1.3\ \mathrm{W\,m^{-2}\,K^{-1}}$, then a direct heating that was permitted to cause a warming of $0.1$ K corresponds to a total of $66$ TW, about 3.6 times present use, and a warming of $0.3$ K to $199$ TW, about eleven times. These figures are indicative, since the response to a heat flux differs from the response to a forcing of another kind, but they show that a planetary energy supply of the scale contemplated by some of the more ambitious proposals, a thousand times the present, could not be used on the surface without a major change in the climate. The conclusion has consequences for the orbital programs of the fifth part: energy generated in space and used on the surface is eventually dissipated on the surface, and it is not a way around the limit; and processes that need very large amounts of energy are better located where their waste heat can be radiated directly to space.

The time scale of the thermal response matters to the interventions of the polar and oceanic kind and to the dynamics of termination. The upper ocean has a large heat capacity, and a layer of one hundred meters has a capacity of about $4\times10^8\ \mathrm{J\,m^{-2}\,K^{-1}}$, so that with $\lambda=1.3$ the time constant of the layer's response to a forcing is about ten years. The deep ocean has a far longer time constant. A forcing that is applied and then removed therefore produces a response that lags it by a decade for the upper layer and by centuries for the deep, and a program that reduces a forcing by a deliberate intervention and then ceases will be followed by a rebound whose rate is set by these time constants and by the rate at which the intervention is withdrawn. The termination analysis required by Chapter 16 for any such program must include the thermal response, and the formalization gives the relevant relation.

Radiation presents a different kind of limit. Several of the proposed systems involve nuclear energy, either as a power source for installations in remote locations or in the form of the refrigeration at the polar margins that appears in the author's earlier proposals, and any such system must be assessed against the physics of shielding and the institutions of safeguards. The attenuation of a narrow beam of gamma radiation by a material of thickness $x$ follows the exponential law with an attenuation coefficient $\mu$, which for photons of one million electron volts in ordinary concrete is about $0.15\ \mathrm{cm^{-1}}$. The thickness that halves the beam is $4.7$ cm, the thickness that reduces it tenfold is $15.7$ cm, and a reduction by a factor of $10^6$ requires about $94$ cm. These are figures for an idealized beam; real shielding is complicated by scattering, by the build-up of secondary radiation, by neutrons that penetrate far more readily, and by activation of the shield itself, but the exponential law shows why a few meters of material suffice for a reactor and why the shielding of a structure that must be occupied is a question of mass and not of principle. The more difficult problems are not the shielding but the fuel cycle, the waste, and the institutions that guard against diversion. The weight of the arguments recorded in Chapter 21 about dual use applies here with equal force, and a nuclear installation in a polar region, where jurisdiction is contested, adds the complication that the body competent to inspect and license it may not exist.

The third constraint, resilience, is the property on which the preceding chapters most depend and which is hardest to demonstrate. An infrastructure is resilient to the extent that it continues to function under shocks, and it is restorable to the extent that it can be returned to function after failing, with limited help from outside. The two properties are different. A system with great redundancy may be robust to the failure of any component and yet impossible to restart if it should fail altogether, because the knowledge, the spare parts, and the energy that a restart requires are held outside it. The author's earlier work on the Caldera Reactor proposes a layered architecture of redundancy, described as hyperpleonastic, with ten recursive layers of safety and fault tolerance. The attraction of such an architecture is evident from the arithmetic of independent failure, since if each of ten layers has a probability of failure of one in ten then the probability that all fail is one in ten billion. The arithmetic is misleading when the layers share causes of failure, a flood, a fabrication defect, a software error, an act of sabotage, that disable many of them at once. With a modest common-cause fraction of five per cent the probability of total failure is not $10^{-10}$ but $5\times10^{-3}$, set almost entirely by the common-cause term, as the formalization shows and as Chapter 29 develops. The lesson is that depth of redundancy cannot substitute for diversity, and that the layers of a defensive architecture should differ in physical principle, in supplier, in location, and in the people who maintain them.

The measure of restorability proposed here is an index of autonomy. For an installation that must restore itself, the relevant stocks are those of the items on which restoration depends: energy, spare parts, consumables, tools, and the hours of trained labor available. Each is consumed at some rate during the restoration, and the time for which the installation can continue is set by the stock that runs out first. The autonomy time is the minimum over the items of the stock divided by the consumption rate, and restoration is feasible without outside help only if the autonomy time exceeds the time required to complete it. The structure is identical to the minimum law of Chapter 19, and it has the same implication: the addition of stock to an item that is not the binding one does not improve autonomy, and the sensible policy is to identify the binding item, increase its stock until another item binds, and repeat. The practical value of the index is that it can be measured by exercises, in which an installation is deprived of outside supply and required to restart, and that the results can be published. The program requires that every installation which claims resilience supply the figure and the record of the exercises on which it rests.

\section*{Formalization}

First, the heat balance. Let the surface temperature anomaly $T$ obey $C\,dT/dt=F(t)-\lambda T$, where $C$ is the heat capacity per unit area of the responding layer, $F$ the forcing (in $\mathrm{W\,m^{-2}}$), and $\lambda>0$ the feedback parameter.

\begin{proposition}
For constant $F$ the solution with $T(0)=0$ is $T(t)=(F/\lambda)(1-e^{-t/\tau})$ with $\tau=C/\lambda$. If $F$ is removed at time $t_1$, then $T(t)=T(t_1)e^{-(t-t_1)/\tau}$ for $t>t_1$. If a standing cooling forcing $-F_c$ has held the system at the equilibrium $T=-F_c/\lambda$ and is removed abruptly, the initial rate of warming is $F_c/C$, and the rate is smaller the more gradually the forcing is withdrawn.
\end{proposition}

\begin{proof}
The equation is linear with constant coefficients and the solutions are standard. For the last statement, immediately after removal $C\,dT/dt=0-\lambda T=F_c$.
\end{proof}

For a mixed layer of $h=100$ m, $C=\rho c_ph=1025\times3990\times100=4.09\times10^8\ \mathrm{J\,m^{-2}\,K^{-1}}$, and with $\lambda=1.3$ the time constant is $\tau=3.15\times10^8$ s $=10.0$ years. An abrupt termination of a cooling forcing of $F_c=1\ \mathrm{W\,m^{-2}}$ implies an initial warming rate of $F_c/C=2.4\times10^{-9}$ K s$^{-1}=0.077$ K per year in the layer, which is the quantity the termination plan must bound by staging the withdrawal over an interval comparable to or longer than $\tau$.

For the energy ceiling, let $A=5.1\times10^{14}\ \mathrm{m^2}$ be the surface area. A heat flux $P$ spread uniformly produces a forcing $P/A$ and an equilibrium warming $P/(\lambda A)$, so that for a tolerated warming $\Delta T$ the maximum sustained power is $P_{\max}=\lambda A\Delta T$. With $\lambda=1.3$ this is $66$ TW for $\Delta T=0.1$ K and $199$ TW for $0.3$ K. The absorbed solar power is $S(1-\alpha)A/4=1.21\times10^{17}$ W for $S=1361\ \mathrm{W\,m^{-2}}$ and albedo $\alpha=0.30$, and the present use of $18.5$ TW is $1.5\times10^{-4}$ of it.

Second, the attenuation law. A narrow beam of intensity $I_0$ traversing a thickness $x$ of material with linear attenuation coefficient $\mu$ is attenuated to $I=I_0e^{-\mu x}$. The half-value layer is $x_{1/2}=\ln2/\mu$ and the tenth-value layer $x_{1/10}=\ln10/\mu$; to achieve an attenuation factor $10^n$ requires thickness $n\ln10/\mu$. With $\mu=0.147\ \mathrm{cm^{-1}}$ (concrete of density $2.3\ \mathrm{g\,cm^{-3}}$ and mass attenuation coefficient $0.064\ \mathrm{cm^2\,g^{-1}}$ at $1$ MeV), $x_{1/2}=4.7$ cm, $x_{1/10}=15.7$ cm, and $n=6$ requires $94$ cm. These are narrow-beam values for primary photons only, and a shield design requires transport calculations that include build-up and neutrons.

Third, the redundancy of layered defenses and the autonomy index. Let a defense have $k$ layers, each failing with probability $p$ and let a common cause with probability $\beta p$ disable all layers, the remaining failures being independent with probability $(1-\beta)p$ each.

\begin{proposition}
The probability that every layer fails is $p_{\mathrm{eff}}=\beta p+(1-\beta p)\big((1-\beta)p\big)^k$, which is bounded below by $\beta p$ for every $k$, and which tends to $\beta p$ as $k\to\infty$ when $(1-\beta)p<1$. Hence for $p=0.1$, $\beta=0.05$ and $k=10$, $p_{\mathrm{eff}}=5.0\times10^{-3}$, although the independent model gives $p^k=10^{-10}$.
\end{proposition}

\begin{proof}
The total failure event is the union of the common-cause event, of probability $\beta p$, and, in its complement, the event that all $k$ independent failures occur, of probability $((1-\beta)p)^k$. Add the probabilities.
\end{proof}

The autonomy time for items $i=1,\dots,n$ with stocks $S_i$ and consumption rates $c_i>0$ is $T_a=\min_iS_i/c_i$, and restoration of duration $T_r$ is feasible without external supply if and only if $T_a\ge T_r$. If $T_r$ is required to be at most the autonomy time, then the cost-effective increment is to raise the stock of the item attaining the minimum until a second item attains it, and the proof is the same as that of the minimum-law proposition of Chapter 19.
